% GENERATED FILE. Edit canonical lessons, then run npm run generate:books.
% canonical-source-hash: fnv1a64:0e2710237ff555f6
% canonical-lessons: FA-C77-panjah, FA-C77-hezar, FA-C77-donya, FA-C77-rang, FA-C77-ahang

\chapter{Fifty, Thousand, World, Colour, Song}
\label{ch:fa-fifty-thousand-world-colour-song}
\hlchaptermodality{77}

\begin{chapteropening}
\textbf{By the end of this chapter:} \emph{I can say fifty, a thousand, the world, a colour and a song, a tune.}

Complete the last lesson of chapter 77: I can say fifty, a thousand, the world, a colour and a song, a tune.
\end{chapteropening}

\section[\texorpdfstring{\fa{پنجاه} (panjâh)}{پنجاه (panjâh)}]{\fa{پنجاه} (panjâh) — fifty}

\label{lesson:FA-C77-panjah}



\begin{quote}
Before the new one: say the Persian for thirty, then the Persian for forty.
\end{quote}

\subsection*{You'll want to know: \fa{پنجاه}}
\textbf{\fa{پنجاه}} — \emph{panjâh} — \textquotedblleft{}fifty\textquotedblright{}.

\begin{cousinweb}
\emph{panjâh} is a native Persian word, inherited rather than borrowed; it is usually traced to Middle Persian \emph{panjāh}, \textquotedblleft{}fifty\textquotedblright{}. It is built on \emph{panj}, \textquotedblleft{}five\textquotedblright{}.
\end{cousinweb}

\begin{grammarlens}[title={What you've built}]
One more word for this chapter.
\end{grammarlens}

\subsection*{Your turn}
\begin{itemize}
\raggedright
  \item \emph{Say it:} \emph{panjâh}
  \item \emph{Say it:} \emph{panjâh}, once more
  \item \emph{Say it:} say \emph{panjâh}
  \item \emph{Recall:} say the Persian for thirty, then the Persian for forty, then say \emph{panjâh} again
\end{itemize}

\begin{tcolorbox}[breakable,colback=teal!4,colframe=teal!35!black,title={Before you move on}]
What does \fa{پنجاه} mean? (\textquotedblleft{}Fifty\textquotedblright{}.) Say it once more.
\end{tcolorbox}
\section[\texorpdfstring{\fa{هزار} (hezâr)}{هزار (hezâr)}]{\fa{هزار} (hezâr) — a thousand}

\label{lesson:FA-C77-hezar}



\begin{quote}
Before the new one: say the Persian for forty, then the Persian for fifty.
\end{quote}

\subsection*{You'll want to know: \fa{هزار}}
\textbf{\fa{هزار}} — \emph{hezâr} — \textquotedblleft{}a thousand\textquotedblright{}.

\begin{cousinweb}
\emph{hezâr} is a native Persian word, inherited rather than borrowed; it is usually traced to Middle Persian \emph{hazār}, \textquotedblleft{}thousand\textquotedblright{}. EtymDB takes it back to Proto-Iranian \emph{*hazahram}, \textquotedblleft{}thousand\textquotedblright{}; Sanskrit \emph{sahasra}, \textquotedblleft{}thousand\textquotedblright{}, is the same Indo-Iranian word.
\end{cousinweb}

\begin{grammarlens}[title={What you've built}]
One more word for this chapter.
\end{grammarlens}

\subsection*{Your turn}
\begin{itemize}
\raggedright
  \item \emph{Say it:} \emph{hezâr}
  \item \emph{Say it:} \emph{hezâr}, once more
  \item \emph{Say it:} say \emph{hezâr}
  \item \emph{Recall:} say the Persian for forty, then the Persian for fifty, then say \emph{hezâr} again
\end{itemize}

\begin{tcolorbox}[breakable,colback=teal!4,colframe=teal!35!black,title={Before you move on}]
What does \fa{هزار} mean? (\textquotedblleft{}A thousand\textquotedblright{}.) Say it once more.
\end{tcolorbox}
\section[\texorpdfstring{\fa{دنیا} (donyâ)}{دنیا (donyâ)}]{\fa{دنیا} (donyâ) — the world}

\label{lesson:FA-C77-donya}



\begin{quote}
Before the new one: say the Persian for fifty, then the Persian for a thousand.
\end{quote}

\subsection*{You'll want to know: \fa{دنیا}}
\textbf{\fa{دنیا}} — \emph{donyâ} — \textquotedblleft{}the world\textquotedblright{}.

\begin{cousinweb}
\emph{donyâ} is borrowed from Arabic \emph{dunyā}, \textquotedblleft{}the nearer world, this life\textquotedblright{}. Its Arabic root is \emph{d-n-w} (Lane). It is the word for \textquotedblleft{}nearer\textquotedblright{}: this world, as against the next.
\end{cousinweb}

\begin{grammarlens}[title={What you've built}]
One more word for this chapter.
\end{grammarlens}

\subsection*{Your turn}
\begin{itemize}
\raggedright
  \item \emph{Say it:} \emph{donyâ}
  \item \emph{Say it:} \emph{donyâ}, once more
  \item \emph{Say it:} say \emph{donyâ}
  \item \emph{Recall:} say the Persian for fifty, then the Persian for a thousand, then say \emph{donyâ} again
\end{itemize}

\begin{tcolorbox}[breakable,colback=teal!4,colframe=teal!35!black,title={Before you move on}]
What does \fa{دنیا} mean? (\textquotedblleft{}The world\textquotedblright{}.) Say it once more.
\end{tcolorbox}
\section[\texorpdfstring{\fa{رنگ} (rang)}{رنگ (rang)}]{\fa{رنگ} (rang) — a colour}

\label{lesson:FA-C77-rang}



\begin{quote}
Before the new one: say the Persian for a thousand, then the Persian for the world.
\end{quote}

\subsection*{You'll want to know: \fa{رنگ}}
\textbf{\fa{رنگ}} — \emph{rang} — \textquotedblleft{}a colour\textquotedblright{}.

\begin{cousinweb}
\emph{rang} is a native Persian word, inherited rather than borrowed: it continues Middle Persian \emph{rang}, \textquotedblleft{}colour, dye\textquotedblright{} (EtymDB).
\end{cousinweb}

\begin{grammarlens}[title={What you've built}]
One more word for this chapter.
\end{grammarlens}

\subsection*{Your turn}
\begin{itemize}
\raggedright
  \item \emph{Say it:} \emph{rang}
  \item \emph{Say it:} \emph{rang}, once more
  \item \emph{Say it:} say \emph{rang}
  \item \emph{Recall:} say the Persian for a thousand, then the Persian for the world, then say \emph{rang} again
\end{itemize}

\begin{tcolorbox}[breakable,colback=teal!4,colframe=teal!35!black,title={Before you move on}]
What does \fa{رنگ} mean? (\textquotedblleft{}A colour\textquotedblright{}.) Say it once more.
\end{tcolorbox}
\section[\texorpdfstring{\fa{آهنگ} (âhang)}{آهنگ (âhang)}]{\fa{آهنگ} (âhang) — a song, a tune}

\label{lesson:FA-C77-ahang}



\begin{quote}
Before the new one: say the Persian for the world, then the Persian for a colour.
\end{quote}

\subsection*{You'll want to know: \fa{آهنگ}}
\textbf{\fa{آهنگ}} — \emph{âhang} — \textquotedblleft{}a song, a tune\textquotedblright{}.

\begin{cousinweb}
\emph{âhang} is a native Persian word, inherited rather than borrowed: it continues Middle Persian \emph{āhang}, \textquotedblleft{}pull, draught; melody\textquotedblright{} (EtymDB).
\end{cousinweb}

\begin{grammarlens}[title={What you've built}]
One more word for this chapter.
\end{grammarlens}

\subsection*{Your turn}
\begin{itemize}
\raggedright
  \item \emph{Say it:} \emph{âhang}
  \item \emph{Say it:} \emph{âhang}, once more
  \item \emph{Say it:} say \emph{âhang}
  \item \emph{Recall:} say the Persian for the world, then the Persian for a colour, then say \emph{âhang} again
\end{itemize}

\begin{tcolorbox}[breakable,colback=teal!4,colframe=teal!35!black,title={Before you move on}]
What does \fa{آهنگ} mean? (\textquotedblleft{}A song, a tune\textquotedblright{}.) Say it once more.
\end{tcolorbox}
